% Общий преамбул для конспекта. Собирается через XeLaTeX/LuaLaTeX.

\usepackage{polyglossia}
\setmainlanguage{russian}
\setotherlanguage{english}

\usepackage{fontspec}
\setmainfont{PT Serif}
% Без Scale=MatchLowercase Menlo визуально заметно крупнее и жирнее PT Serif
% при том же кегле — инлайн-код "прыгает" на фоне обычного текста.
% MatchLowercase подгоняет размер моноширинного шрифта по x-height основного.
\setmonofont{Menlo}[Scale=MatchLowercase]
\newfontfamily\cyrillicfont{PT Serif}
\newfontfamily\cyrillicfonttt{Menlo}[Scale=MatchLowercase]

\usepackage{amsmath, amssymb, amsthm, mathtools}

% amsthm по умолчанию продолжает текст доказательства на той же строке, что
% и заголовок "Доказательство." — здесь после заголовка всегда перенос
% строки. Копия \newenvironment{proof} из amsthm.sty с добавленным \newline.
\makeatletter
\renewenvironment{proof}[1][\proofname]{\par
  \pushQED{\qed}%
  \normalfont \topsep6\p@\@plus6\p@\relax
  \trivlist
  \item[\hskip\labelsep
        \itshape
    #1\@addpunct{.}]\hspace{0pt}\\\ignorespaces
}{%
  \popQED\endtrivlist\@endpefalse
}
\makeatother
\usepackage{microtype}
\usepackage{enumitem}
\usepackage{graphicx}
\graphicspath{{images/}}
\usepackage{xcolor}

% Документ читается с экрана, не печатается — поля уже, чем в книжной
% вёрстке, и одинаковые с обеих сторон (документ не переплетается).
\usepackage[margin=2.7cm]{geometry}

% Чуть больше межстрочного интервала, чем "книжный" одинарный — легче
% воспринимать длинные доказательства и оставлять место для мысленных пометок.
\linespread{1.08}

% Приглушённый серый — для служебных пометок (ссылка на Lean-исходник),
% чтобы они не конкурировали по контрасту с математическим текстом.
\definecolor{leangray}{gray}{0.58}

% Приглушённый синий — для кликабельных ссылок на другие аксиомы/теоремы
% (\taoref). hyperref подключён с [hidelinks], то есть сам он ссылки никак
% не выделяет — цвет здесь единственный сигнал, что по тексту можно
% кликнуть, поэтому красим текст ссылки вручную сами.
\definecolor{taoblue}{RGB}{37,84,150}

% Компактные, но не слипшиеся списки по умолчанию — пригодится для шагов
% доказательства/условий упражнений.
\setlist{itemsep=0.2em, topsep=0.5em}

\usepackage[hidelinks]{hyperref}
\usepackage{xurl} % разрешает перенос \path/\url по любому символу, не только после "_"/"."
\usepackage{cleveref}

% hyperref по умолчанию делает \url/\path курсивным, если курсивен окружающий
% текст (\urlstyle{same}) — для моноширинного шрифта это выглядит плохо
% (синтетический наклон). Код должен быть прямым всегда.
\urlstyle{tt}

% Заголовки разделов вида "3.2" — как у Тао
\renewcommand\thesection{\thechapter.\arabic{section}}

% Внутренние (немые) amsthm-окружения. Их автосчётчик не используется по
% значению — он лишь механизм amsthm для печати номера. Сам номер мы каждый
% раз подставляем вручную через \renewcommand\theaxiomEnv, см. ниже.
%
% Везде используется стиль "definition" (жирный заголовок, ПРЯМОЙ текст тела) —
% стандартный "plain" даёт курсивное тело, и курсив на длинных абзацах и
% особенно в моноширинном шрифте читается плохо.
\theoremstyle{definition}
\newtheorem{axiomEnv}{Аксиома}
\newtheorem{theoremEnv}[axiomEnv]{Теорема}
\newtheorem{propositionEnv}[axiomEnv]{Утверждение}
\newtheorem{lemmaEnv}[axiomEnv]{Лемма}
\newtheorem{corollaryEnv}[axiomEnv]{Следствие}
\newtheorem{definitionEnv}[axiomEnv]{Определение}
\newtheorem{exampleEnv}[axiomEnv]{Пример}
\newtheorem{exerciseEnv}[axiomEnv]{Упражнение}
\newtheorem{assumptionEnv}[axiomEnv]{Предположение}
\newtheorem*{axiomEnv*}{Аксиома}
\newtheorem*{theoremEnv*}{Теорема}
\newtheorem*{propositionEnv*}{Утверждение}
\newtheorem*{lemmaEnv*}{Лемма}
\newtheorem*{corollaryEnv*}{Следствие}
\newtheorem*{definitionEnv*}{Определение}
\newtheorem*{exampleEnv*}{Пример}
\newtheorem*{exerciseEnv*}{Упражнение}
\newtheorem*{assumptionEnv*}{Предположение}
\newtheorem*{remark}{Замечание}

% Дословная нумерация Тао. Использование:
%   \begin{axiom}{3.8}{универсальная спецификация} ... \end{axiom}
% Первый аргумент — номер БУКВАЛЬНО как в докстринге Lean-теоремы (сверяться
% через \leanref ниже), второй — заголовок в скобках (можно оставить {}).
% У Тао аксиомы нумеруются "Глава.N" (Axiom 3.8), а теоремы/леммы/... —
% "Глава.Раздел.N" (Proposition 3.4.9), поэтому общая авто-нумерация по
% разделу здесь не годится — номер всегда указывается вручную.
% Каждое окружение заодно ставит \label{префикс:номер} — это даёт возможность
% ссылаться на него настоящей кликабельной ссылкой через \taoref (см. ниже),
% без ручной возни с метками. Коллизий не бывает: номера Тао уникальны по
% всей книге.
%
% \nopagebreak в конце запрещает разрыв страницы сразу после формулировки —
% без него между "Утверждение 2.1.6." и следующим за ним "Доказательство."
% может влезть граница страницы, что рвёт нить рассуждения.
\newenvironment{axiom}[2]{\renewcommand\theaxiomEnv{#1}\begin{axiomEnv}[#2]\label{ax:#1}}{\end{axiomEnv}\nopagebreak}
\newenvironment{theorem}[2]{\renewcommand\theaxiomEnv{#1}\begin{theoremEnv}[#2]\label{thm:#1}}{\end{theoremEnv}\nopagebreak}
\newenvironment{proposition}[2]{\renewcommand\theaxiomEnv{#1}\begin{propositionEnv}[#2]\label{prop:#1}}{\end{propositionEnv}\nopagebreak}
\newenvironment{lemma}[2]{\renewcommand\theaxiomEnv{#1}\begin{lemmaEnv}[#2]\label{lem:#1}}{\end{lemmaEnv}\nopagebreak}
\newenvironment{corollary}[2]{\renewcommand\theaxiomEnv{#1}\begin{corollaryEnv}[#2]\label{cor:#1}}{\end{corollaryEnv}\nopagebreak}
\newenvironment{definition}[2]{\renewcommand\theaxiomEnv{#1}\begin{definitionEnv}[#2]\label{def:#1}}{\end{definitionEnv}\nopagebreak}
\newenvironment{example}[2]{\renewcommand\theaxiomEnv{#1}\begin{exampleEnv}[#2]\label{ex:#1}}{\end{exampleEnv}\nopagebreak}
\newenvironment{exercise}[2]{\renewcommand\theaxiomEnv{#1}\begin{exerciseEnv}[#2]\label{exr:#1}}{\end{exerciseEnv}\nopagebreak}
\newenvironment{assumption}[2]{\renewcommand\theaxiomEnv{#1}\begin{assumptionEnv}[#2]\label{assum:#1}}{\end{assumptionEnv}\nopagebreak}

% Беззвёздная версия — для утверждений, которые у Тао не пронумерованы
% (например, приведены как рассуждение в прозе, без формальной метки).
% Аргумент — заголовок в скобках, номер не печатается.
\newenvironment{axiom*}[1]{\begin{axiomEnv*}[#1]}{\end{axiomEnv*}\nopagebreak}
\newenvironment{theorem*}[1]{\begin{theoremEnv*}[#1]}{\end{theoremEnv*}\nopagebreak}
\newenvironment{proposition*}[1]{\begin{propositionEnv*}[#1]}{\end{propositionEnv*}\nopagebreak}
\newenvironment{lemma*}[1]{\begin{lemmaEnv*}[#1]}{\end{lemmaEnv*}\nopagebreak}
\newenvironment{corollary*}[1]{\begin{corollaryEnv*}[#1]}{\end{corollaryEnv*}\nopagebreak}
\newenvironment{definition*}[1]{\begin{definitionEnv*}[#1]}{\end{definitionEnv*}\nopagebreak}
\newenvironment{example*}[1]{\begin{exampleEnv*}[#1]}{\end{exampleEnv*}\nopagebreak}
\newenvironment{exercise*}[1]{\begin{exerciseEnv*}[#1]}{\end{exerciseEnv*}\nopagebreak}
\newenvironment{assumption*}[1]{\begin{assumptionEnv*}[#1]}{\end{assumptionEnv*}\nopagebreak}

% Ссылка на конкретную лемму/теорему в Lean-файле того же раздела.
% \path (из url.sty, с xurl) даёт моноширинный текст, переносимый по любому
% символу — длинные идентификаторы не вылезают за поля, и не нужно
% экранировать подчёркивания вручную.
\newcommand{\leanref}[1]{\path{#1}}

% Короткая вставка кода прямо в тексте (не длинный идентификатор — для тех
% есть \leanref). \mbox запрещает перенос строки внутри — короткий фрагмент
% вроде "zero : Nat" не должен разрываться пополам на границе строки.
\newcommand{\code}[1]{\mbox{\texttt{#1}}}

% Вынесенная строка с именем Lean-теоремы/леммы вместе с её формальной
% сигнатурой — ставится последней внутри окружения, после
% доказательства/формулировки. #1 — идентификатор, #2 — явные параметры и
% формулировка (без неявных {..}-параметров), в родном синтаксисе Lean.
% Флеш-лефт, а не hfill: сигнатура обычно длиннее одного идентификатора и
% переносится на несколько строк, вправо это смотрелось бы плохо.
\newcommand{\leansource}[2]{\par\nobreak\smallskip\noindent{\footnotesize\color{leangray}Lean:~{\leanref{#1}}\ \texttt{#2}}\par}

% Единообразный разбор случаев внутри \begin{proof}...\end{proof}:
%   \caselabel{$\Omega \in \Omega$} Тогда ...
\newcommand{\caselabel}[1]{\medskip\noindent\textbf{Случай #1.}\ }

% То же самое, но с произвольной подписью — для шагов индукции и похожих
% структурных частей доказательства:
%   \prooflabel{База} $n = 0$: ...
%   \prooflabel{Переход} ...
\newcommand{\prooflabel}[1]{\medskip\noindent\textbf{#1.}\ }

% Кликабельная ссылка на другую аксиому/теорему/... по её метке (см. префиксы
% выше — ax:, thm:, prop:, lem:, cor:, def:, ex:, exr:, assum:). Падеж и
% формулировку текста задаём сами вторым аргументом, макрос не гадает.
%   По \taoref{ax:2.1}{аксиоме 2.1}, $0$ — натуральное число.
\newcommand{\taoref}[2]{\hyperref[#1]{\color{taoblue}#2}}

% Обоснование шага доказательства — мелкой серой пометкой сразу после
% содержательного утверждения, вместо ведущего "По аксиоме X, ...". Само
% рассуждение остаётся обычным текстом на переднем плане, ссылка на
% аксиому/лемму отходит на второй план (как и \leansource):
%   $0$ — натуральное число.\by{\taoref{ax:2.1}{аксиома 2.1}}
\newcommand{\by}[1]{\ {\footnotesize\color{leangray}(#1)}}
